\section{The Merge-Split Algorithm}
\label{sec:algorithm}

\subsection{Algorithm Specification}

\begin{definition}[Wilson UST]
\label{def:ust}
$\mathrm{Wilson\_UST}(G, \mathrm{rng})$ denotes Wilson's loop-erased random
walk algorithm~\citep{wilson1996} that produces a uniformly random spanning
tree of graph $G$ using the random source $\mathrm{rng}$.
Two calls with independent random sources produce independent spanning trees.
Expected time: $O(|E(G)| \log |E(G)|)$.
\end{definition}

The \mergesplit{} chain uses two deterministic seed streams per step,
derived from the base seed and step index as follows:
\begin{align}
  \mathrm{step\_seed}(b, t)
    &= \mathrm{SHA256}\!\left(
         \texttt{"MS\_STEP\_"} \,\|\, t^{\mathrm{le64}} \,\|\,
         \texttt{"\_"} \,\|\, 0^{\mathrm{le32}} \,\|\,
         \texttt{"\_"} \,\|\, b^{\mathrm{le64}}
       \right)[0{:}8]
       \;\text{as } u64, \label{eq:step-seed}\\
  \mathrm{rev\_seed}(b, t)
    &= \mathrm{SHA256}\!\left(
         \texttt{"MS\_REVERSE\_"} \,\|\, t^{\mathrm{le64}} \,\|\,
         \texttt{"\_"} \,\|\, 0^{\mathrm{le32}} \,\|\,
         \texttt{"\_"} \,\|\, b^{\mathrm{le64}}
       \right)[0{:}8]
       \;\text{as } u64, \label{eq:rev-seed}
\end{align}
where $b$ is the base seed, $t$ is the step index, $\cdot^{\mathrm{le64}}$
and $\cdot^{\mathrm{le32}}$ denote little-endian 8-byte and 4-byte encodings
respectively, and $[0{:}8]$ takes the first 8 bytes of the SHA-256 output.

\begin{algorithm}
\caption{\mergesplit{} Chain}
\label{alg:mergesplit}
\begin{algorithmic}[1]
\Require Initial plan $\pi_0$, steps $S$, base seed $b$
\Ensure  Selected plan $\pi^*$ (minimum EC over visited plans)
\State $\pi \gets \pi_0$
\State $\mathrm{best} \gets (\ec(\pi_0),\, \pi_0)$
\For{$t \gets 1$ \textbf{to} $S$}
  \State $\mathrm{rng}_{\mathrm{fwd}} \gets \mathrm{SmallRng::seed\_from\_u64}(\mathrm{step\_seed}(b, t))$
  \State $\mathrm{rng}_{\mathrm{rev}} \gets \mathrm{SmallRng::seed\_from\_u64}(\mathrm{rev\_seed}(b, t))$
  \State $(d_i, d_j) \gets \mathrm{sample\_adjacent\_pair}(\pi,\; \mathrm{rng}_{\mathrm{fwd}})$
         \Comment{with pair reselection on failure}
  \State $R \gets \pi^{-1}(d_i) \cup \pi^{-1}(d_j)$
  \State $T_{\mathrm{fwd}} \gets \mathrm{Wilson\_UST}(G[R],\; \mathrm{rng}_{\mathrm{fwd}})$
  \State $\mathrm{cuts\_fwd} \gets |\{\text{balanced cuts of } T_{\mathrm{fwd}}\}|$
  \If{$\mathrm{cuts\_fwd} = 0$} \textbf{continue} \EndIf
  \State $e^* \gets \mathrm{rng}_{\mathrm{fwd}}.\mathrm{choose}(\text{balanced cuts of } T_{\mathrm{fwd}})$
  \State $\pi' \gets$ plan induced by removing $e^*$ from $T_{\mathrm{fwd}}$
  \State $T_{\mathrm{rev}} \gets \mathrm{Wilson\_UST}(G[R],\; \mathrm{rng}_{\mathrm{rev}})$
         \Comment{independent sample}
  \State $\mathrm{cuts\_rev} \gets |\{\text{balanced cuts of } T_{\mathrm{rev}}\}|$
  \If{$\mathrm{cuts\_rev} = 0$}
    \State $\alpha \gets 0.0$
  \Else
    \State $\alpha \gets \mathrm{cuts\_fwd} / \mathrm{cuts\_rev}$
         \Comment{$= Q(\pi' \!\to\! \pi) / Q(\pi \!\to\! \pi')$ in expectation}
  \EndIf
  \If{$\mathrm{rng}_{\mathrm{fwd}}.\mathrm{sample}(\mathrm{Uniform}[0,1]) < \min(1, \alpha)$}
    \State $\pi \gets \pi'$
    \If{$\ec(\pi') < \mathrm{best.ec}$}
      \State $\mathrm{best} \gets (\ec(\pi'),\, \pi')$
    \EndIf
  \EndIf
\EndFor
\State \Return $\mathrm{best.plan}$
\end{algorithmic}
\end{algorithm}

\paragraph{Design choices.}
Three aspects of Algorithm~\ref{alg:mergesplit} deserve comment.

First, the forward and reverse random sources ($\mathrm{rng}_{\mathrm{fwd}}$
and $\mathrm{rng}_{\mathrm{rev}}$) are derived from different SHA-256 prefixes
(\texttt{"MS\_STEP\_"} vs.\ \texttt{"MS\_REVERSE\_"}), ensuring they are
statistically independent.
Using the same source for both trees would reintroduce the symmetry collapse
of standard \recom{}: if $T_{\mathrm{rev}}$ is derived from $T_{\mathrm{fwd}}$,
the cut counts would be correlated and the ratio would no longer estimate
$Q(\pi' \!\to\! \pi) / Q(\pi \!\to\! \pi')$.

Second, if $\mathrm{cuts\_rev} = 0$, the ratio is set to $0.0$ and the
proposal is rejected.
This is necessary for well-definedness: if $T_{\mathrm{rev}}$ has no balanced
cut, the reverse move from $\pi'$ back to $\pi$ has probability zero under
any spanning tree of $G[R]$ that has no balanced cuts.
Accepting such a proposal would violate detailed balance.

Third, the chain tracks the best plan (minimum edge cut) across all visited
plans, rather than the full visited sequence.
This is appropriate for the plan-selection use case; for full ensemble
analysis, the visited sequence should be retained.

\paragraph{CLI invocation.}
\begin{verbatim}
redist state --state NC --year 2020 \
    --search merge-split --merge-split-steps 1000 --percentile 0.0
\end{verbatim}

\subsection{The Acceptance Ratio: Direction Derivation}

\begin{proposition}[Ratio direction]
\label{prop:ratio-direction}
The Metropolis-Hastings acceptance ratio for the move $\pi \to \pi'$
under the uniform target distribution is:
\[
  \frac{Q(\pi' \to \pi)}{Q(\pi \to \pi')}
  = \frac{1 / |\mathcal{C}(T_{\mathrm{rev}}, R)|}
         {1 / |\mathcal{C}(T_{\mathrm{fwd}}, R)|}
  = \frac{|\mathcal{C}(T_{\mathrm{fwd}}, R)|}{|\mathcal{C}(T_{\mathrm{rev}}, R)|}
  = \frac{\mathrm{cuts\_fwd}}{\mathrm{cuts\_rev}}.
\]
The forward cut count appears in the \textbf{numerator} (not the denominator).
\end{proposition}

\begin{proof}
The proposal probability for the forward move $\pi \to \pi'$ is:
\[
  Q(\pi \to \pi')
  = \underbrace{\frac{1}{|\mathcal{P}(\pi)|}}_{%
      \text{pair selection}}
  \cdot \underbrace{\frac{1}{|\mathcal{T}(R)|}}_{%
      \text{UST sample}}
  \cdot \underbrace{\frac{1}{|\mathcal{C}(T_{\mathrm{fwd}}, R)|}}_{%
      \text{cut selection}}.
\]
The forward cut count $|\mathcal{C}(T_{\mathrm{fwd}}, R)|$ appears in the
\emph{denominator} of $Q(\pi \to \pi')$: more balanced cuts means each
individual cut is selected with lower probability.

The proposal probability for the reverse move $\pi' \to \pi$ is:
\[
  Q(\pi' \to \pi)
  = \frac{1}{|\mathcal{P}(\pi')|}
  \cdot \frac{1}{|\mathcal{T}(R)|}
  \cdot \frac{1}{|\mathcal{C}(T_{\mathrm{rev}}, R)|}.
\]

Assuming $|\mathcal{P}(\pi)| = |\mathcal{P}(\pi')|$ (the number of
adjacent district pairs is the same in both plans, which holds when
the adjacency structure is preserved), the ratio is:
\[
  \frac{Q(\pi' \to \pi)}{Q(\pi \to \pi')}
  = \frac{|\mathcal{C}(T_{\mathrm{fwd}}, R)|}{|\mathcal{C}(T_{\mathrm{rev}}, R)|}.
\]
The implementation correctly uses
$\mathrm{cuts\_fwd} / \mathrm{cuts\_rev}$. \qed
\end{proof}

\begin{remark}[Intuition for the direction]
A large forward cut count means the forward proposal probability
$Q(\pi \to \pi')$ is diluted (many cuts competed for selection, and this
particular cut won by chance).
If the reverse tree $T_{\mathrm{rev}}$ has a smaller cut count, the reverse
move $\pi' \to \pi$ is less diluted and thus more probable.
A higher-probability reverse move means $\pi'$ tends to ``want'' to return
to $\pi$ more strongly, which raises the acceptance probability.
Hence: $\mathrm{cuts\_fwd} / \mathrm{cuts\_rev} > 1$ leads to elevated
acceptance; $\mathrm{cuts\_fwd} / \mathrm{cuts\_rev} < 1$ leads to reduced
acceptance.
\end{remark}

\subsection{Approximate Detailed Balance}

\begin{theorem}[Approximate detailed balance]
\label{thm:approx-balance}
The \mergesplit{} chain satisfies detailed balance in expectation with
respect to the uniform distribution over valid redistricting plans.
That is, for any two valid plans $\pi$ and $\pi'$ that differ by a
single region merge-split move:
\[
  \mathbb{E}\!\left[
    \mathrm{Trans}(\pi, \pi')
  \right]
  = \mathbb{E}\!\left[
    \mathrm{Trans}(\pi', \pi)
  \right],
\]
where $\mathrm{Trans}(\pi, \pi')$ is the transition probability from $\pi$
to $\pi'$ and the expectation is over the randomness of the spanning tree
samples.
\end{theorem}

\begin{proof}[Proof sketch]
By symmetry of Wilson's algorithm, the expected number of balanced cuts
of a uniformly random spanning tree of $G[R]$ is the same regardless of
which of the two plans $\pi$ or $\pi'$ is current.
Therefore:
\[
  \mathbb{E}\!\left[\frac{\mathrm{cuts\_fwd}}{\mathrm{cuts\_rev}}\right]
  = \mathbb{E}\!\left[\frac{|\mathcal{C}(T, R)|}{|\mathcal{C}(T', R)|}\right]
  \approx 1,
\]
where the approximation holds when $|\mathcal{C}(T, R)|$ and
$|\mathcal{C}(T', R)|$ are approximately identically distributed
(both are uniform samples from $\mathcal{T}(R)$).
Under the Metropolis-Hastings rule with ratio $\alpha = \mathrm{cuts\_fwd}/\mathrm{cuts\_rev}$,
the expected transition satisfies detailed balance with the uniform
distribution as the target.
The exact argument follows \citet{janson2022}, Proposition~2.
\end{proof}

\subsection{Propositions}

\begin{proposition}[Positive acceptance rate]
\label{prop:positive-accept}
On a connected census-tract graph with at least one valid two-district
split in each region, the acceptance rate of \mergesplit{} is strictly
positive.
\end{proposition}

\begin{proof}
If $G[R]$ is connected and at least one balanced cut exists in some
spanning tree of $G[R]$, then $\mathrm{cuts\_fwd} > 0$ with positive
probability (over the choice of $T_{\mathrm{fwd}}$).
Similarly, $\mathrm{cuts\_rev} > 0$ with positive probability.
When both counts are positive, the acceptance probability is
$\min(1, \mathrm{cuts\_fwd}/\mathrm{cuts\_rev}) > 0$. \qed
\end{proof}

\begin{proposition}[Per-step cost]
\label{prop:cost}
Each step of \mergesplit{} requires $O(m \log m)$ expected time, where
$m = |E(G[R])|$ is the number of edges in the merged region.
\end{proposition}

\begin{proof}
The dominant operations are the two calls to Wilson's algorithm, each
$O(m \log m)$ expected time~\citep{wilson1996}.
Enumerating balanced cuts is $O(m)$ (a linear pass over the tree edges).
All other operations (pair selection, plan update) are $O(k)$ or $O(1)$.
Total: $O(m \log m)$. \qed
\end{proof}
